\section{Technical Supplement: Mathematical Foundations of OPH}\label{technical-supplement-mathematical-foundations-of-observer-patch-holography}

\begin{quote}
The quantitative branches discussed in this paper use two external inputs (\(P\) for pixel area and \(N_{\mathrm{scr}}\) for screen capacity) together with structural axioms and explicit regulator/scaling premises. This supplement collects the mathematical ingredients for the conditional Lorentz/null-modular/Einstein branch, conditional compact gauge reconstruction, and the related bosonic quantitative sectors.
\end{quote}

\subsection{For Theoretical Physicists}\label{for-theoretical-physicists}

This document collects the mathematical derivations and stated premises used in the OPH paper suite. The Lorentz/null-modular/Einstein and compact-gauge statements are scaling-limit / categorical results under explicit regularity hypotheses; these hypotheses are properties of the EFT phase or UV completion under discussion, not fixed-cutoff matrix-algebra theorems.

\textbf{Scope convention.} Statements about Lorentz recovery, modular flow, and the Einstein branch are scaling-limit / EFT statements rather than literal fixed-cutoff type-I claims. Statements about gauge reconstruction refer to the refinement-stable directed colimit of fixed-cutoff sector categories, not to the finite block list at one regulator scale.

\begin{center}\rule{0.5\linewidth}{0.5pt}\end{center}

\section{Reference Point}\label{1-core-axioms-and-definitions}

The five axioms, the canonical technical-premise ledger \(T1\)--\(T7\), and the dependency map appear in the main OPH theorem ledger. This supplement restates the fixed-cutoff realized presentation and the theorem-local derived controls used repeatedly below. MAR enters only in the gauge-selection portions; the Lorentz/null-modular and Einstein discussions mainly use the first four axioms plus the technical premises stated where needed.

\subsection{Technical Premises and Fixed-Cutoff Realization}\label{12-regulator-premises}

The supplement follows the canonical premise scheme of the main manuscript. The only labeled external premises are \(T1\)--\(T7\). The fixed-cutoff realized presentation, the local-Gibbs form, quasi-local propagation, endpoint-Lipschitz interval control, and refinement-stable branch persistence described below are unpackings or consequences of the OPH axioms rather than additional labeled assumptions.

\textbf{Fixed-cutoff realized presentation.} On a finite regulator cellulation, each regulator cell \(x\) carries a finite-dimensional factor \(\mathfrak h_x\), and local patch algebras are finite type-I algebras. Gauge-as-gluing is realized as a compact boundary redundancy action on cut data, so physical algebras are fixed-point algebras of the corresponding boundary action. This is the realized regulator form of Axiom 1 plus overlap consistency; it is not a separate premise layer.

\textbf{Local finite-constraint MaxEnt branch.} Axiom 3 states that at scale \(\ell_{\mathrm{UV}}\) the realized state maximizes entropy subject to a finite set of gauge-invariant local constraints
\[
\mathcal C_{\ell_{\mathrm{UV}}}=\{O_a(x)\}_{a=1}^{N_{\mathrm{con}}},
\]
with each \(O_a(x)\) supported in a ball of radius \(O(\ell_{\mathrm{UV}})\) and the label set \(a\) independent of the number of regulator cells. On a finite-dimensional algebra, the standard Lagrange-multiplier argument therefore gives
\[
\omega_{\ell_{\mathrm{UV}}}(\lambda)
=
Z_{\ell_{\mathrm{UV}}}(\lambda)^{-1}
\exp\!\bigl(-K_{\ell_{\mathrm{UV}}}(\lambda)\bigr),
\]
with
\[
K_{\ell_{\mathrm{UV}}}(\lambda)
=
\sum_x \sum_{a=1}^{N_{\mathrm{con}}}\lambda_a O_a(x)
+\sum_b \mu_b Q_b,
\]
where the \(Q_b\) are the optional global conserved-charge constraints. Thus the logarithm of the selected state is itself a quasi-local generator built from the same bounded-support densities that define the branch.

\textbf{Derived quasi-local propagation.} Because \(K_{\ell_{\mathrm{UV}}}(\lambda)\) is a bounded-support local sum on a finite regulator graph, standard Lieb--Robinson theory applies to the automorphism group it generates, or to any branch generator lying in the same bounded-support algebraic closure. Writing \(d(X,Y)\) for the regulator graph distance between local supports \(X\) and \(Y\), one obtains constants \(C,\xi,v_{\ell_{\mathrm{UV}}}<\infty\) such that
\[
\bigl\|[\tau_t^{\ell_{\mathrm{UV}}}(A_X),B_Y]\bigr\|
\le
C\,\|A_X\|\,\|B_Y\|\,\min(|X|,|Y|)\,
e^{-(d(X,Y)-v_{\ell_{\mathrm{UV}}}|t|)/\xi}.
\]
This is the fixed-cutoff propagation control later used in the collar and null-strip discussions. It is not a separate technical premise beyond Axiom 3 plus the finite regulator presentation.

\textbf{Derived endpoint control.} The same bounded-support structure gives the only endpoint regularity used later. Whenever a later interval- or collar-labeled fixed-cutoff generator on this branch is obtained by restricting the same local density list to the chosen region and then separating off the central endpoint term, changing the region from \(I\) to \(I'\) adds or removes only an \(O(|I'\Delta I|)\) collar of local terms. Hence on any fixed local-energy-bounded domain,
\[
\left|\langle\psi,(K[I']-K[I])\phi\rangle\right|
\le
C_{\psi,\phi,I_{\max}}\,|I'\Delta I|.
\]
This is the endpoint-Lipschitz / finite-variation control imported later into the null bridge. No second independent regulator premise is needed.

\textbf{Internal refinement branch.} Axiom 3 also states that the same finite local constraint family is preserved under refinement. Therefore the realized states at different cutoffs lie in one common finite-dimensional multiplier space \(\lambda\). For any compatible family of refinement channels \(\Phi_{\ell\to L}\), one may write on the realized branch
\[
\Phi_{\ell\to L}\bigl(\omega_\ell(\lambda)\bigr)
=
\omega_L\!\bigl(R_{\ell\to L}(\lambda)\bigr)
\]
for an induced map \(R_{\ell\to L}\) on that multiplier space. The phrase ``refinement-stable branch'' therefore means persistence along a trajectory or invariant subset of this finite-dimensional map. This state-side refinement notion is enough to compare realized states across cutoffs. It does \emph{not} by itself upgrade fixed-cutoff edge labels to transportable refinement-persistent sectors, and it does not show that the directed colimit of such sectors is nontrivial. In particular, when the gauge branch later speaks of a refinement-stable directed colimit of sectors, the entire transportable directed-system/fiber clause is the explicit content of \(T7\); Axiom 3 supplies only the realized state branch along which one asks whether such sector transport persists.

\textbf{What remains external.} This internalization claim stops at the fixed-cutoff local-Gibbs/Lieb--Robinson/refinement package. Quantitative exponential mixing or exact-Markov replacement, when invoked later, is separate state information tied to the recoverability side of Axiom 4 rather than to Axiom 3 alone. The genuinely external technical premises are therefore exactly the canonical list \(T1\)--\(T7\): the Bisognano--Wichmann geometric scaling branch and its scope clause, fixed-cap stationarity, the null-stress/density-upgrade burden, vanishing transport obstruction when global sector transportability is invoked, and the transportable sector-colimit plus symmetric/bosonic/fiber-functor conditions used in gauge reconstruction. No separate regulator premise for quasi-local propagation or refinement stability survives once Axiom 3 is unpacked in this way.

\textbf{Proposition 1.0a (Ancilla-stable UV underdetermination).} Let a fixed-cutoff OPH realization be stabilized by finite ancillary factors \(K_P\) in a fixed product state, with observable patch algebras embedded as \(\mathcal A(P)\otimes \mathbf 1_{K_P}\) and repair dynamics acting trivially on the ancillas. Then observable expectations on the physical subalgebras, overlap data, the local-Gibbs branch, the collar conditional mutual information \(I(A:D\mid B)\), the Fawzi--Renner remainder, the collar Markov modulus, and the quotient normal form are unchanged. Thus OPH determines the UV branch only modulo such ancillary stabilization together with gauge or implementation hiding, not a unique microscopic presentation.

\textbf{Proof.} Product ancillas leave physical observables unchanged, cancel additively inside conditional mutual information, and are inert under the repair maps. Hence every invariant listed above is unchanged. QED.

\subsection{Edge-Center Completion (EC)}\label{13-edge-center-completion-ec}

At fixed regulator scale, overlap consistency itself determines the boundary redundancy. On the
ordinary or central-defect branch, EC is the invariant-theoretic consequence of that derived
boundary action. Exact Markovity is a further state property and is not implied by EC alone.

\textbf{Proposition 1.1a (Derived boundary gluing datum).} Choose a finite regulator chart for the patches meeting along a connected cut \(\Sigma\). Because the local overlap algebras are finite-dimensional matrix algebras, any overlap-consistent recharting is an inner automorphism and is implemented by a unitary on the cut Hilbert space. The compact closure of the subgroup generated by these recharting unitaries is a compact boundary redundancy group \(K_\Sigma\). If triple-overlap defects are central, the projective composition law lifts to a compact central extension \(\widehat K_\Sigma\); on the ordinary branch one simply sets \(\widehat K_\Sigma = K_\Sigma\). A genuinely noncentral 2-group defect is the only obstruction to reducing the overlap transition system to an ordinary compact group action, and it is handled below by the higher-gauge replacement rather than left as a residual fixed-cutoff gap.

\textbf{Proof.} On a finite-dimensional matrix algebra every \(^*\)-automorphism is inner, so each overlap-consistent recharting is implemented by a unitary on the cut Hilbert space. The subgroup generated by those unitaries has compact closure inside a finite-dimensional unitary group. A central defect gives a projective composition law that lifts to a central extension; a genuinely noncentral defect is handled by the higher-gauge replacement stated below. QED.

\textbf{Theorem 1.1 (Derived EC decomposition).} Under the fixed-cutoff regulator realization of Section 1.2, and on the ordinary or central-defect branch, the collar Hilbert space is
\[
\mathcal H_{B_\delta}
=
(\tilde{\mathcal H}_{B_L}\otimes \tilde{\mathcal H}_{B_R})^{\widehat K_\Sigma}
\cong
\bigoplus_{\alpha}
\left(\mathcal H_{b_L^\alpha}\otimes \mathcal H_{b_R^\alpha}\right),
\]
and the center of the collar algebra is generated by the block projectors:
\[
Z(\mathcal A(B_\delta))
=
\bigoplus_\alpha \mathbb C\cdot \mathbf 1_\alpha.
\]
The right half-collar carries the contragredient representation because it sees the inverse transport across the same cut.

\textbf{Proof.} Decompose the left boundary data into irreps
\[
\tilde{\mathcal H}_{B_L}
=
\bigoplus_\alpha
\left(W_\alpha\otimes \mathcal H_{b_L^\alpha}\right),
\]
and the right boundary data into dual irreps
\[
\tilde{\mathcal H}_{B_R}
=
\bigoplus_\beta
\left(W_\beta^*\otimes \mathcal H_{b_R^\beta}\right).
\]
Then
\[
\tilde{\mathcal H}_{B_L}\otimes \tilde{\mathcal H}_{B_R}
=
\bigoplus_{\alpha,\beta}
(W_\alpha\otimes W_\beta^*)\otimes
(\mathcal H_{b_L^\alpha}\otimes \mathcal H_{b_R^\beta}).
\]
By Schur's lemma,
\[
(W_\alpha\otimes W_\beta^*)^{\widehat K_\Sigma}
\cong
\begin{cases}
\mathbb C, & \alpha=\beta,\\
0, & \alpha\neq \beta,
\end{cases}
\]
so only matched left/right labels survive, and the surviving block projectors generate the center. QED.

\textbf{Proposition 1.1b (Derived higher-gauge cut datum).} On the genuinely noncentral branch, weak overlap gluing on a connected cut \(\Sigma\) is encoded by a compact crossed module \(\mathbb K_\Sigma=(H_\Sigma\xrightarrow{\partial_\Sigma}G_\Sigma,\triangleright)\) with defect class
\[
q_\Sigma\in \check H^2(N_\Sigma,H_\Sigma\to G_\Sigma),
\]
and compact higher-gauge change system
\[
\mathcal T_\Sigma=C^1(N_\Sigma,H_\Sigma)\rtimes C^0(N_\Sigma,G_\Sigma).
\]
\textbf{Proof.} Pair-overlap rechartings are inner; triple-overlap associators promote the transition system to a compact unitary \(2\)-group; and skeletal strictification identifies that \(2\)-group with crossed-module data. Because the nerve is finite, the higher-gauge change system is compact. QED.

\textbf{Theorem 1.1c (Higher-gauge EC decomposition and defect transport).} On the genuinely noncentral branch,
\[
\mathcal H_{B_\delta}^{2g}
=
(\tilde{\mathcal H}_{B_L}\otimes \tilde{\mathcal H}_{B_R})^{\mathcal T_\Sigma}
\cong
\bigoplus_\lambda
(\mathcal H_{b_L^\lambda}\otimes \mathcal H_{b_R^\lambda}),
\]
and
\[
Z(\mathcal A_{2g}(B_\delta))
=
\bigoplus_\lambda \mathbb C\cdot \mathbf 1_\lambda.
\]
Moreover the higher-gauge defect class \(q_\Sigma\) is invariant under local rechartings, vanishes iff the defect is removable, and classifies fixed-cutoff genuinely noncentral sectors.
\textbf{Proof.} Finite-dimensional unitary \(\mathcal T_\Sigma\)-modules decompose into
irreducibles, and Schur matching leaves only equal left/right labels inside the invariant tensor
product. The transport statement is the crossed-module \v{C}ech analogue of the central-defect
case: coboundaries change representatives while preserving classes, the trivial class strictifies
the cocycle, and a nontrivial class cannot be gauged away. QED.

\textbf{Corollary 1.2 (Exact Markov adds the state factorization).} On either the ordinary/central branch of Theorem 1.1 or the genuinely noncentral higher-gauge branch of Theorem 1.1c, if, in addition,
\[
I_\omega(A_\delta:D_\delta\mid B_\delta)=0,
\]
or one passes to the explicitly stated idealized recoverability limit that reduces to exact Markovity, then
\[
\rho_{A_\delta B_\delta D_\delta}
=
\bigoplus_\alpha p_\alpha
\left(\rho_{A_\delta b_L^\alpha}\otimes \rho_{b_R^\alpha D_\delta}\right).
\]
EC therefore gives the kinematic block decomposition; exact Markovity is the extra state input that gives the HJPW normal form.

\textbf{Remark.} Approximate recoverability gives controlled deviations from this normal form; it is not implied by EC alone. The fixed-cutoff topological UV package is closed on the ordinary, central-defect, and genuinely noncentral higher-gauge branches. The later open burden is the controlled refinement or scaling regime, together with the transportable-sector, modular, or null-limit premises used by the refinement-limit gravity and compact-gauge branches.
\subsection{External Continuous Inputs}\label{14-fundamental-parameters}

In the quantitative branches discussed in this paper, OPH uses exactly two external continuous inputs:

\begin{enumerate}
\def\labelenumi{\arabic{enumi}.}
\tightlist
\item
 \textbf{Pixel area:} \(P \equiv a_{\text{cell}}/\ell_P^2 \approx 1.63\)
\item
 \textbf{Screen capacity:} \(N_{\mathrm{scr}} \equiv \log(\dim \mathcal{H}_{\text{tot}}) \sim 10^{122}\)
\end{enumerate}

From these:

\begin{itemize}
\tightlist
\item
 Newton\textquotesingle s constant: \(G = \frac{a_{\text{cell}}}{4\bar{\ell}(t)}\)
\item
 Cosmological constant: \(\Lambda = \frac{3\pi}{G N_{\mathrm{scr}}}\)
\end{itemize}

\begin{center}\rule{0.5\linewidth}{0.5pt}\end{center}

\section{Emergence of Gravity and Quantum Mechanics}\label{2-emergence-of-gravity-and-quantum-mechanics}

\subsection{Local Lorentz Structure}\label{21-local-lorentz-structure}

\textbf{Theorem 2.1 (Conformal Group Isomorphism):} The orientation-preserving conformal group of \(S^2\) is isomorphic to the connected Lorentz group: \[\text{Conf}^+(S^2) \cong \text{PSL}(2,\mathbb{C}) \cong \text{SO}^+(3,1)\]

This provides the kinematic structure of special relativity.

\subsection{2.2 Geometric Modular Flow}\label{22-geometric-modular-flow}

At fixed cutoff, the collar analysis gives a literal density matrix on each cap and confines the modular defect to a shrinking collar up to explicit carried errors. The Lorentz claim is therefore not a finite-matrix identity. Its target is the scaling-limit cap pair
\[
(\mathcal A_\infty(C),\omega_\infty^C)
\]
selected by the controlled refinement/scaling clause \textbf{T3}. No independent Euclidean-regularity or cap-isotropy premise is used here.

For a shrinking collar family \((A_\delta,B_\delta,D_\delta)\) around \(\partial C\), write
\[
\varepsilon_\delta:=I(A_\delta:D_\delta\mid B_\delta)_\omega,
\qquad
r_{\mathrm{FR}}(\varepsilon_\delta):=
2\sqrt{1-e^{-\varepsilon_\delta}}
\le 2\sqrt{\varepsilon_\delta},
\]
and, on each fixed faithful collar model,
\[
\eta_\delta^{\mathrm M}
:=
4\lambda_\ast^{-1}\,
\delta^{\mathrm M}_{A_\delta:B_\delta:D_\delta}(\varepsilon_\delta).
\]
All exact collar formulas in this subsection are therefore to be read either literally at exact Markovity or along a controlled collar family satisfying
\[
\delta^{\mathrm M}_{A_\delta:B_\delta:D_\delta}(\varepsilon_\delta)\to 0
\qquad(\delta\downarrow 0),
\]
with \(r_{\mathrm{FR}}(\varepsilon_\delta)\) and \(\eta_\delta^{\mathrm M}\) carried explicitly.

For each cap \(C\), let \(\lambda_C(s)\subset \mathrm{Conf}^+(S^2)\) denote the standard cap-preserving conformal one-parameter subgroup, normalized so that the null blow-up near a smooth cut acts by \(v\mapsto e^{-s}v\), and let \(\alpha_{\lambda_C(s)}\) denote the induced automorphism of the scaling-limit cap net.

\textbf{Theorem 2.2 (Geometric modular flow on caps on the OPH Bisognano--Wichmann branch).} Assume the OPH axioms, the derived fixed-cutoff collar package, the controlled scaling-limit scope clause \textbf{T3}, and the branch condition \textbf{T1} that the realized refinement-stable scaling branch for caps and their conformal images lies in the Bisognano--Wichmann geometric modular class. Then the scaling-limit modular automorphism group of the cap pair is
\[
\sigma_t^{\omega_\infty^C}
=
\alpha_{\lambda_C(2\pi t)}.
\]
At finite regulator stage the nonadditive part of
\[
K_C^{(\delta)}:=-\log \rho_C^{(\delta)}
\]
is confined to the shrinking collar up to the carried errors \(r_{\mathrm{FR}}(\varepsilon_\delta)\) and \(\eta_\delta^{\mathrm M}\). If the limit cap algebra remains type I, this automorphism identity is represented by
\[
K_C=2\pi B_C;
\]
otherwise the automorphism identity is the full statement and the action is outer.

\textbf{Proof.} Markov locality does one precise job: it reduces the modular problem to the shrinking collar, with the finite-stage deviation from an exact Markov reference carried explicitly by \(r_{\mathrm{FR}}(\varepsilon_\delta)\) and \(4\lambda_\ast^{-1}\delta^{\mathrm M}_{A_\delta:B_\delta:D_\delta}(\varepsilon_\delta)\). By \textbf{T3}, the theorem concerns only the scaling-limit cap pair. By \textbf{T1}, on the selected branch the scaling-limit modular automorphism group is the standard cap-preserving conformal subgroup up to normalization:
\[
\sigma_t^{\omega_\infty^C}
=
\alpha_{\lambda_C(\kappa_C t)}
\]
for some \(\kappa_C>0\). The modular group is KMS by definition, and because \(s\) is normalized by unit null dilation on the branch, the Bisognano--Wichmann relation fixes \(\kappa_C=2\pi\). Thus
\[
\sigma_t^{\omega_\infty^C}
=
\alpha_{\lambda_C(2\pi t)}.
\]
If the limit algebra remains type I, this automorphism identity can be represented by \(K_C=2\pi B_C\); otherwise the automorphism identity is the full statement and the action is outer. QED.

\textbf{BW boundary and remaining UV scaffold.} This theorem proves the BW conclusion on the \textbf{T1}-selected branch and fixes its \(2\pi\) normalization without a Euclidean-regularity selector. It does not prove that MaxEnt alone selects that branch. The remaining UV-side scaffold is exactly:
\begin{enumerate}[leftmargin=2em]
\item a realized transported cap-local system, consisting of the cap-local test family, the projectively compatible transported marginal family, and the asymptotic transport-equivalence certificate, together with the derived carried-collar schedule
\[
r_{\mathrm{FR}}(\varepsilon_{n,m,\delta})
+
4\bar\lambda_{m,\delta}^{-1}
\delta^{\mathrm M}_{A_{n,m,\delta}:B_{n,m,\delta}:D_{n,m,\delta}}
\bigl(\varepsilon_{n,m,\delta}\bigr),
\]
whose only nonlatent faithfulness-side input is the eventual fixed-local-collar modular-transport common floor on the transported physical marginals; and
\item ordered cut-pair rigidity on the emitted scaling-limit cap pair.
\end{enumerate}
No second comparison-state faithfulness clause is missing: once the exact-Markov modulus tends to zero on a fixed collar model, the exact-Markov comparison marginals inherit the same eventual common floor.

\begin{corollary}[Lorentz kinematics on the BW branch]\label{cor:lorentz}
Under the hypotheses of Theorem~\ref{thm:bw},
\[
\mathrm{Conf}^+(S^2)\cong \mathrm{PSL}(2,\mathbb C)\cong \mathrm{SO}^+(3,1).
\]
Hence the cap modular flows are the standard Lorentz boost/dilation subgroups in the celestial-sphere realization, and the induced local kinematic group is the connected Lorentz group.
\end{corollary}

\begin{proof}
Orientation-preserving conformal maps of \(S^2\) are exactly the M\"obius transformations, so
\[
\mathrm{Conf}^+(S^2)\cong \mathrm{PSL}(2,\mathbb C)\cong \mathrm{SO}^+(3,1).
\]
By Theorem~\ref{thm:bw}, on the \textbf{T1} branch the realized scaling-limit cap modular automorphism groups are the standard cap-preserving conformal dilations with the \(2\pi\) normalization fixed internally. Hence the local kinematics induced by modular flow is the connected Lorentz group.
\end{proof}
\subsection{Null Modular Bridge at Fixed Cutoff}\label{23-null-modular-bridge-eft-conditions-n1-n3}

The fixed-cutoff null bridge reaches the derived half-sided-inclusion step. What is established here is: transferred cut-center data on null strips, exact-or-controlled strip additivity on one inherited strip model, endpoint-Lipschitz control of the renormalized half-line family, and the geometric blow-up argument that turns the null half-line net into a half-sided modular pair, so that Borchers--Wiesbrock yields an explicit positive null-translation generator on its Stone domain~\cite{wiesbrock93}. What is not established here is the downstream density upgrade or the relativistic null-stress identification. This theorem-local package should therefore be read with the following split.

\textbf{Fixed-cutoff null-strip bridge.} For consecutive null strips
\[
I_-=(v_1,v_2),\qquad
J=(v_2,v_3),\qquad
I_+=(v_3,v_4),
\]
assume the two cuts inherit the same compact overlap-redundancy data as ordinary collars, so that on a compatible type-I regulator presentation
\[
\tilde{\mathcal H}_{J}
\cong
\bigoplus_{\alpha_-,\alpha_+}
W_{\alpha_-}^{*}\otimes
\mathcal H_{j^{\alpha_-,\alpha_+}}\otimes
W_{\alpha_+}
\]
and the strip algebra is the commutant of the transported cut action. Then
\[
\mathcal A(J)
\cong
\bigoplus_{\alpha_-,\alpha_+}
\mathcal B(\mathcal H_{j^{\alpha_-,\alpha_+}}),
\qquad
Z(\mathcal A(J))
=
\bigoplus_{\alpha_-,\alpha_+}\mathbb C\,P_{\alpha_-,\alpha_+}.
\]
If, in addition, each multiplicity space admits an inherited split
\[
\mathcal H_{j^{\alpha_-,\alpha_+}}
\cong
\mathcal H_{j_L^{\alpha_-,\alpha_+}}
\otimes
\mathcal H_{j_R^{\alpha_-,\alpha_+}},
\]
with \(\mathcal A(I_-\cup J)\) acting blockwise only on
\(\mathcal H_{i_-^{\alpha_-}}\otimes \mathcal H_{j_L^{\alpha_-,\alpha_+}}\)
and \(\mathcal A(J\cup I_+)\) acting blockwise only on
\(\mathcal H_{j_R^{\alpha_-,\alpha_+}}\otimes \mathcal H_{i_+^{\alpha_+}}\),
then exact Markov strip states satisfy exact four-term modular additivity on the canonical inherited model:
\[
\Delta K_J(\omega)
:=
K_{I_-\cup J\cup I_+}(\omega)-K_{I_-\cup J}(\omega)-K_{J\cup I_+}(\omega)+K_J(\omega)
=0,
\]
equivalently up to a central endpoint term. For a general strip state with
\[
I(I_-:I_+\mid J)_\omega\le \varepsilon,
\]
define
\[
\mathfrak M_{I_-:J:I_+}
:=
\left\{
\sigma:\ I(I_-:I_+\mid J)_\sigma=0
\right\},
\qquad
\delta^{\mathrm M}_{I_-:J:I_+}(\varepsilon)
:=
\sup\left\{
\inf_{\sigma\in\mathfrak M_{I_-:J:I_+}}\|\rho-\sigma\|_1:
I(I_-:I_+\mid J)_\rho\le\varepsilon
\right\}.
\]
Choose \(\widetilde\omega_J\in \mathfrak M_{I_-:J:I_+}\) with
\[
\|\omega-\widetilde\omega_J\|_1
\le
\delta^{\mathrm M}_{I_-:J:I_+}(\varepsilon),
\]
and set
\[
\mathfrak D_J(\omega,\widetilde\omega_J)
:=
\Delta K_J(\omega)-\Delta K_J(\widetilde\omega_J).
\]
Then
\[
K_{I_-\cup J\cup I_+}(\omega)
=
K_{I_-\cup J}(\omega)+K_{J\cup I_+}(\omega)-K_J(\omega)
+K_{\partial,J}(\widetilde\omega_J)
+\mathfrak D_J(\omega,\widetilde\omega_J),
\]
with \(K_{\partial,J}(\widetilde\omega_J)\) central. If the relevant marginals are uniformly faithful with lower spectral bound \(\lambda_\ast>0\), then
\[
\|\mathfrak D_J(\omega,\widetilde\omega_J)\|_\infty
\le
4\lambda_\ast^{-1}\,
\delta^{\mathrm M}_{I_-:J:I_+}(\varepsilon).
\]
Independently, Fawzi--Renner recovery~\cite{fawziRenner15} supplies a constructive comparison state with
\[
r_{\mathrm{FR}}(\varepsilon)
=
2\sqrt{1-e^{-\varepsilon}}
\le
2\sqrt{\varepsilon}.
\]
Thus this null-strip bridge is exact only at exact Markovity, or else along a controlled strip family on one fixed inherited strip model with
\[
\delta^{\mathrm M}_{I_-:J:I_+}(\varepsilon_J)\to 0.
\]

\textbf{Derived half-sided modular inclusion on the null half-line blow-up net.} Fix a smooth cut point, choose affine coordinate \(v\) on the chosen null generator \(\Omega\) so that the cut sits at \(v=0\), and define for \(a\ge 0\)
\[
H_a:=(a,\infty),
\qquad
\mathcal M_a(\Omega)
:=
\overline{\bigvee_{a<c<d<\infty}\mathcal A((c,d),\Omega)}.
\]
By isotony of the null interval net,
\[
a\le b
\quad\Longrightarrow\quad
\mathcal M_b(\Omega)\subseteq \mathcal M_a(\Omega).
\]
On the OPH geometric scaling branch, the blow-up of cap modular flow acts by the null dilation
\[
v-v_0\mapsto e^{-2\pi t}(v-v_0).
\]
Hence
\[
\sigma_t^\omega\bigl(\mathcal M_a(\Omega)\bigr)
=
\mathcal M_{e^{-2\pi t}a}(\Omega).
\]
For \(t\le 0\) one has \(e^{-2\pi t}a\ge a\), so \(H_{e^{-2\pi t}a}\subseteq H_a\) and therefore
\[
\sigma_t^\omega\bigl(\mathcal M_a(\Omega)\bigr)\subseteq \mathcal M_a(\Omega).
\]
Thus \(\mathcal M_a(\Omega)\subset \mathcal M_0(\Omega)\) is half-sided modular. After the harmless convention change \(t\mapsto -t\), this is the standard positive-time form. The half-sided modular inclusion is therefore derived from OPH null-interval structure, isotony, and the scaling-limit geometric action rather than imported separately.

\textbf{Endpoint-Lipschitz control and weak tail generator.} For renormalized modular Hamiltonians
\[
\widetilde K[I,\Omega]
:=
K[I,\Omega]-K_\partial(I,\Omega),
\qquad
\widetilde K_a(\Omega):=\widetilde K[(a,\infty),\Omega],
\]
bounded intervals satisfy
\[
\bigl|
\langle\psi,(\widetilde K[(a',b'),\Omega]-\widetilde K[(a,b),\Omega])\phi\rangle
\bigr|
\le
C_{\psi,\phi,\Omega}\bigl(|a'-a|+|b'-b|\bigr),
\]
and half-lines satisfy
\[
\bigl|
\langle\psi,(\widetilde K_{a'}(\Omega)-\widetilde K_a(\Omega))\phi\rangle
\bigr|
\le
C_{\psi,\phi,\Omega}|a'-a|.
\]
Hence, for fixed \(\psi,\phi\),
\[
f_{\psi,\phi}(a):=\langle\psi,\widetilde K_a(\Omega)\phi\rangle
\]
is locally Lipschitz and absolutely continuous, and its distributional derivative defines the weak tail generator
\[
\langle\psi,q(a,\Omega)\phi\rangle
:=
-\frac{1}{2\pi}\,\partial_a f_{\psi,\phi}(a).
\]
For \(a<b\),
\[
\langle\psi,(\widetilde K_b(\Omega)-\widetilde K_a(\Omega))\phi\rangle
=
-2\pi\int_a^b \langle\psi,q(v,\Omega)\phi\rangle\,dv.
\]
This is the end of the fixed-cutoff analytic part of the bridge. The stronger assumptions
\[
q \ \text{weakly differentiable},
\qquad
f_{\psi,\phi}(a)\to 0,
\qquad
q_{\psi,\phi}(a)\to 0
\quad (a\to+\infty)
\]
would permit the downstream density formula
\[
\langle\psi,p(a,\Omega)\phi\rangle
:=
-\partial_a\langle\psi,q(a,\Omega)\phi\rangle
=
\frac{1}{2\pi}\,\partial_a^2 f_{\psi,\phi}(a),
\qquad
\langle\psi,\widetilde K_a(\Omega)\phi\rangle
=
2\pi\int_a^\infty (v-a)\,\langle\psi,p(v,\Omega)\phi\rangle\,dv,
\]
but that density upgrade is not established here.

\textbf{Imported Borchers consequence.} Let \(\Delta_0(\Omega)\) be the modular operator of the standard pair \((\mathcal M_0(\Omega),\omega)\), define \(K_0(\Omega):=-\log \Delta_0(\Omega)\), and let \(\Delta_a(\Omega)\) denote the modular operator of \((\mathcal M_a(\Omega),\omega)\). Applying Borchers--Wiesbrock to the derived half-sided pair above~\cite{wiesbrock93} yields a unique strongly continuous unitary group \(U_\Omega(a)=e^{iaP_\Omega}\) with positive self-adjoint Stone generator \(P_\Omega\) on
\[
D(P_\Omega)
:=
\left\{
\psi\in\mathcal H:
\lim_{a\to0}\frac{U_\Omega(a)\psi-\psi}{ia}\ \text{exists}
\right\}.
\]
It obeys
\[
\Delta_0(\Omega)^{it}U_\Omega(a)\Delta_0(\Omega)^{-it}
=
U_\Omega(e^{-2\pi t}a),
\qquad
[K_0(\Omega),P_\Omega]=-\,i\,2\pi P_\Omega,
\]
and if \(K_a(\Omega):=-\log \Delta_a(\Omega)\), then
\[
K_a(\Omega)=K_0(\Omega)-2\pi a\,P_\Omega
\]
as a quadratic-form identity on \(D(K_0(\Omega))\cap D(P_\Omega)\).

\textbf{Consequence.} This theorem-local package proves the fixed-cutoff strip bridge through the weak tail generator \(q(a,\Omega)\), derives half-sided modular inclusion on the null half-line blow-up net, and therefore identifies an explicit positive null-translation generator \(P_\Omega\) together with its Stone domain and half-line modular-Hamiltonian relation. The remaining downstream burdens are the density-upgrade hypotheses and the relativistic EFT null-stress premise. Those later steps are explicit rather than hidden inside an additional labeled premise package.
\subsection{Stress Tensor Reconstruction}\label{24-stress-tensor-reconstruction}

\textbf{Lemma 2.3 (Null Data Ambiguity):} If a symmetric tensor \(X_{ab}\) satisfies \(X_{ab}k^a k^b = 0\) for all null \(k\), then: \[X_{ab} = \phi \, g_{ab}\] for some scalar \(\phi\).

\emph{Proof:} Decompose \(X_{ab} = Y_{ab} + \phi g_{ab}\) with \(Y\) traceless. For null \(k = (1, \hat{n})\): \[X_{kk} = Y_{kk} + \phi g_{kk} = Y_{kk}\] since \(g_{kk} = 0\). Expanding: \[Y_{kk} = Y_{00} + 2\hat{n}^i Y_{0i} + \hat{n}^i \hat{n}^j Y_{ij}\]

If this vanishes for all \(\hat{n} \in S^2\), each angular moment (scalar, vector, traceless tensor) vanishes independently, so \(Y_{ab} = 0\). \(\square\)

\textbf{Corollary 2.4:} Null modular data determine \(T_{ab}\) only up to \(\phi g_{ab}\). Consequently, the Einstein equation is fixed only up to \(\Lambda g_{ab}\).

\textbf{Reconstruction Formulas:} From \(f(\hat{n}) := T_{kk}(\hat{n})\): \[T_{0i} = \frac{3}{2}\langle \hat{n}_i f \rangle_{S^2}\] \[T_{ij} - \frac{\delta_{ij}}{3}T^{(\text{spatial})}_{kk} = \frac{15}{2}\left\langle \left(\hat{n}_i \hat{n}_j - \frac{\delta_{ij}}{3}\right) f \right\rangle_{S^2}\]

\subsection{Entanglement Equilibrium and Einstein\textquotesingle s Equation}\label{25-entanglement-equilibrium-and-einsteins-equation}

\textbf{Theorem 2.5 (Conditional Jacobson-type Derivation):} In the same local Lorentzian scaling regime, under MaxEnt state selection with generalized entropy stationarity, the null-stress identification premise of the relativistic EFT branch, and admissible first variations about a maximally symmetric reference state~\cite{jacobson95,jacobson16},
\[\delta S_{\text{gen}} = 0\]

implies the rest-frame first-variation relation
\[
\delta\!\left(G_{00} + \Lambda g_{00}\right) = 8\pi G \,\delta\langle T_{00} \rangle.
\]

If this first-variation relation holds for all local directions and reference states in the scaling regime, overlap consistency upgrades it to the full tensor equation
\[
G_{ab} + \Lambda g_{ab} = 8\pi G \langle T_{ab} \rangle
\]
modulo the expected \(\Lambda g_{ab}\) ambiguity.

\textbf{Derivation sketch:}

\begin{enumerate}
\def\labelenumi{\arabic{enumi}.}
\item
 Raychaudhuri equation relates area change to \(R_{kk}\): \[\frac{d\theta}{d\lambda} = -\frac{1}{2}\theta^2 - \sigma^2 - R_{kk}\]
\item
 Recoverable Generalized Entropy implies: \[\frac{d}{d\lambda}\left(\frac{A}{4G} + S_{\text{bulk}}\right) \leq 0\]
\item
 Stationarity at MaxEnt gives: \[R_{kk} = 8\pi G \langle T_{kk} \rangle\]
\item
 By Lemma 2.3, this fixes Einstein\textquotesingle s equation up to \(\Lambda g_{ab}\).
\end{enumerate}

\subsection{Newton\textquotesingle s Constant from Edge Entropy}\label{26-newtons-constant-from-edge-entropy}

\textbf{Theorem 2.6:} The gravitational coupling is fixed by the edge entropy density: \[G = \frac{a_{\text{cell}}}{4\bar{\ell}(t)}\]

where \(\bar{\ell}(t) = a_{\text{cell}}/(4\ell_P^2)\) is the dimensionless entropy per cell.

With \(a_{\text{cell}} \approx 1.63 \ell_P^2\): \[\bar{\ell} = \frac{1.63}{4} \approx 0.4077\]

\begin{center}\rule{0.5\linewidth}{0.5pt}\end{center}

\section{The Measurement Problem}\label{3-the-measurement-problem}

\subsection{Records as Central Variables}\label{31-records-as-central-variables}

\textbf{Definition 3.1:} A \emph{record} in patch \(P\) is a family of orthogonal projectors \(\{\Pi_r\} \subset \mathcal{A}(P)\) such that:

\begin{enumerate}
\def\labelenumi{\arabic{enumi}.}
\tightlist
\item
 \textbf{Classicality:} \(\mathcal{R} := \text{vN}(\{\Pi_r\})\) is approximately abelian
\item
 \textbf{Stability:} \(\{\Pi_r\}\) are stable under effective dynamics
\item
 \textbf{Shareability:} \(\Pi_r\) lies in the center of overlap algebras
\end{enumerate}

\subsection{Markov Structure Theorem}\label{32-markov-structure-theorem}

\textbf{Theorem 3.1 (Quantum Markov chain structure on a fixed EC collar).} If \(I(A:C\mid B)=0\) for a tripartite state \(\rho_{ABC}\), then
\[
\mathcal H_B
\cong
\bigoplus_j
\left(
\mathcal H_{b_L^{(j)}}\otimes \mathcal H_{b_R^{(j)}}
\right)
\]
and the state decomposes as
\[
\rho_{ABC}
=
\bigoplus_j
q_j\,
\rho^{(j)}_{A b_L^{(j)}}\otimes \rho^{(j)}_{b_R^{(j)} C}.
\]
The label \(j\) is a classical variable living in the center of \(\mathcal A(B)\).

On the fixed finite-dimensional collar Hilbert space supplied by EC, let
\[
\mathfrak M_{A:B:C}
:=
\left\{
\sigma_{ABC}: I(A:C\mid B)_\sigma=0
\right\}
\]
and define the exact-Markov distance modulus
\[
\delta^{\mathrm M}_{A:B:C}(\varepsilon)
:=
\sup\left\{
\inf_{\sigma\in\mathfrak M_{A:B:C}}\|\rho-\sigma\|_1:
I(A:C\mid B)_\rho\le\varepsilon
\right\}.
\]
Because the state space is compact and conditional mutual information is continuous in finite dimension,
\[
\delta^{\mathrm M}_{A:B:C}(\varepsilon)\longrightarrow0
\qquad
(\varepsilon\downarrow0).
\]

Approximate collars therefore converge to the exact HJPW normal form in the following sense:
either the state is exactly Markov, or one works on one fixed finite-dimensional collar, or on a
compatible pulled-back family, and sends \(\varepsilon\to0\). The supplement does \emph{not}
claim a universal dimension-free one-shot trace-norm bound from small \(I(A:C\mid B)\) directly
to an exact Markov state.

\subsection{Approximate Markov and Recovery}\label{33-approximate-markov-and-recovery}

\textbf{Theorem 3.2 (Fawzi--Renner recovery~\cite{fawziRenner15}).} If
\[
I(A:C\mid B)_\rho\le\varepsilon,
\]
then there exists a CPTP map \(\mathcal R:\mathcal A(B)\to\mathcal A(BC)\) such that the recovered comparison state
\[
\rho^{\mathrm{rec}}_{ABC}
:=
(\mathrm{id}_A\otimes\mathcal R)(\rho_{AB})
\]
obeys
\[
\|\rho_{ABC}-\rho^{\mathrm{rec}}_{ABC}\|_1
\le
2\sqrt{1-e^{-\varepsilon}}
\le
2\sqrt{\varepsilon}.
\]
If \(\varepsilon\) is quoted in bits rather than nats, this bound is equivalently written as \(2\sqrt{\ln 2\,\varepsilon}\).

The recovered state is a constructive comparison state. It need not be exact Markov. To connect approximate recoverability to the exact splice or additivity formulas used later, one must combine Theorem 3.2 with the fixed-collar modulus of Section 3.2. Concretely, if \(\sigma_\varepsilon\in\mathfrak M_{A:B:C}\) is chosen so that
\[
\|\rho-\sigma_\varepsilon\|_1
\le
\delta^{\mathrm M}_{A:B:C}(\varepsilon),
\]
then:
\begin{enumerate}
\item every bounded observable supported on the left side of the collar differs from its exact-Markov splice value by at most \(\|X\|_\infty\,\delta^{\mathrm M}_{A:B:C}(\varepsilon)\);
\item if the relevant marginals are uniformly faithful with lower spectral bound \(\lambda_\ast>0\), then the corresponding modular-additivity defect differs from the exact Markov value by at most
\[
4\lambda_\ast^{-1}\,
\delta^{\mathrm M}_{A:B:C}(\varepsilon).
\]
\end{enumerate}

Thus the supplement carries two distinct errors:
\[
r_{\mathrm{FR}}(\varepsilon)=O(\varepsilon^{1/2})
\]
for constructive one-shot recovery, and
\[
\delta^{\mathrm M}_{A:B:C}(\varepsilon)\to0
\]
for convergence to the exact Markov normal form on a fixed collar. This is the controlled bridge from approximate recoverability to the exact splice and modular-additivity statements used by the later gravity and consensus arguments.

\subsection{Definiteness of Outcomes}\label{34-definiteness-of-outcomes}

\textbf{Theorem 3.3 (Collapse from EC Structure):} Under EC, the post-measurement physical state is exactly a convex mixture over sector labels \(\alpha\): \[\rho = \bigoplus_\alpha p_\alpha \, \rho^{(\alpha)}\]

For any observer whose accessible observables lie in either side algebra:

\begin{enumerate}
\def\labelenumi{\arabic{enumi}.}
\tightlist
\item
 There are \textbf{no interference observables} between different \(\alpha\)-blocks
\item
 Each observer-instance has a definite \(\alpha\) in the classical sense
\end{enumerate}

\emph{Proof:} Cross-block operators \(|i\rangle\langle j|\) for \(i \in \alpha, j \in \alpha'\) (\(\alpha \neq \alpha'\)) are not gauge-invariant, hence not in the physical algebra. \(\square\)

\subsection{Born Rule as Unique Probability Measure}\label{35-born-rule-as-unique-probability-measure}

\textbf{Theorem 3.4 (Gleason):} For \(\dim \mathcal{H} \geq 3\), any measure \(\mu\) on projectors satisfying:

\begin{enumerate}
\def\labelenumi{\arabic{enumi}.}
\tightlist
\item
 Noncontextuality: \(\mu(P)\) depends only on \(P\)
\item
 Additivity: \(\mu(P \vee Q) = \mu(P) + \mu(Q)\) for \(P \perp Q\)
\item
 Normalization: \(\mu(\mathbf{1}) = 1\)
\end{enumerate}

must have the form: \[\mu(P) = \text{Tr}(\rho P)\] for some density operator \(\rho\).

\textbf{Corollary 3.5:} In OPH with the fixed-cutoff type-I regulator presentation, Born\textquotesingle s rule is the unique overlap-consistent probability assignment.

\subsection{Collapse as Conditioning}\label{36-collapse-as-conditioning}

\textbf{Proposition 3.6 (L\"{u}ders Update):} For record projectors \(\{\Pi_r\}\) and state \(\omega\): \[p_r = \omega(\Pi_r), \qquad \omega_r(X) = \frac{\omega(\Pi_r X \Pi_r)}{\omega(\Pi_r)}\]

Then for all \(X\) commuting with \(\mathcal{R}\): \[\omega(X) = \sum_r p_r \, \omega_r(X)\]

This is classical conditional probability on the record algebra.

\begin{center}\rule{0.5\linewidth}{0.5pt}\end{center}

\section{The Cosmological Principle}\label{4-the-cosmological-principle}

\subsection{MaxEnt Symmetry Inheritance}\label{41-maxent-symmetry-inheritance}

\textbf{Lemma 4.1:} Let \(G\) act on the screen by automorphisms \(\alpha_g\). If:

\begin{enumerate}
\def\labelenumi{\arabic{enumi}.}
\tightlist
\item
 The constraint set is \(G\)-invariant
\item
 Von Neumann entropy is \(G\)-invariant: \(S(\alpha_g(\rho)) = S(\rho)\)
\end{enumerate}

Then the MaxEnt optimizer \(\omega\) can be taken \(G\)-invariant. If unique, it is automatically \(G\)-invariant.

\emph{Proof:} If \(\omega\) is a maximizer, so is \(\omega \circ \alpha_{g^{-1}}\) by constraint invariance and entropy invariance. By uniqueness, \(\omega = \omega \circ \alpha_{g^{-1}}\). \(\square\)

\subsection{Isotropy from SO(3)-Invariant Constraints}\label{42-isotropy-from-so3-invariant-constraints}

\textbf{SO(3)-invariant constraints.} Constraints are SO(3)-invariant on \(S^2\).

\textbf{Theorem 4.2:} Under SO(3)-invariant constraints and MaxEnt uniqueness, the reference state \(\omega\) satisfies: \[\omega \circ \alpha_g = \omega \quad \forall g \in \text{SO}(3)\]

The stress tensor has perfect-fluid form: \[\langle T_{ab} \rangle = \rho \, u_a u_b + p(g_{ab} + u_a u_b)\]

\emph{Proof of perfect-fluid form:} Under SO(3), rank-2 tensors decompose into scalar + vector + traceless tensor irreps. Only scalars are invariant. Hence \(\langle T_{0i} \rangle = 0\) and \(\langle T_{ij} \rangle \propto \delta_{ij}\). \(\square\)

\subsection{Homogeneity from Isotropy Everywhere}\label{43-homogeneity-from-isotropy-everywhere}

\textbf{Theorem 4.3 (Schur-type):} Let \((\Sigma, h_{ij})\) be a 3-dimensional Riemannian manifold. If at every point \(p \in \Sigma\), the curvature tensor is SO(3)-invariant (pointwise isotropy), then: \[R_{ijkl}(p) = K(p)(h_{ik}h_{jl} - h_{il}h_{jk})\]

and by the second Bianchi identity, \(\nabla_m K = 0\), so \(K = \text{const}\).

\textbf{Corollary 4.4:} If OPH supplies isotropy for all observers, spatial geometry is a constant-curvature space form: \(S^3\), \(\mathbb{R}^3\), or \(H^3\).

\subsection{FLRW Emergence}\label{44-flrw-emergence}

\textbf{Theorem 4.5:} With:

\begin{enumerate}
\def\labelenumi{\arabic{enumi}.}
\tightlist
\item
 Semiclassical Einstein equation from recoverable generalized entropy / entanglement equilibrium
\item
 Perfect-fluid stress tensor from SO(3)-invariant MaxEnt
\item
 Positive \(\Lambda\) from finite screen capacity
\end{enumerate}

The emergent geometry is FLRW: \[ds^2 = -dt^2 + a(t)^2 \left(\frac{dr^2}{1-kr^2} + r^2 d\Omega^2\right)\]

with Friedmann equations: \[H^2 + \frac{k}{a^2} = \frac{8\pi G}{3}\rho + \frac{\Lambda}{3}\] \[\dot{H} - \frac{k}{a^2} = -4\pi G(\rho + p)\]

\begin{center}\rule{0.5\linewidth}{0.5pt}\end{center}

\section{Supersymmetry and Gauge Coupling Unification}\label{5-supersymmetry-and-gauge-coupling-unification}

\subsection{One-Loop Unification Algebra}\label{51-one-loop-unification-algebra}

At one loop, if couplings unify at \((M_U, \alpha_U)\): \[\alpha_i^{-1}(M_Z) = \alpha_U^{-1} + \frac{b_i}{2\pi}\ln\frac{M_U}{M_Z}\]

Define \(A_i := \alpha_i^{-1}(M_Z)\), \(L := \ln(M_U/M_Z)\). Eliminating \(\alpha_U\): \[L = \frac{2\pi}{b_1 - b_2}(A_1 - A_2)\]

The consistency relation for \(\alpha_s\) (a consistency check when all three couplings calibrate $P$; a genuine prediction in the two-input mode): \[A_3^{\text{pred}} = \frac{b_3 - b_2}{b_1 - b_2}A_1 + \frac{b_1 - b_3}{b_1 - b_2}A_2\]

\subsection{SM vs MSSM Coefficients}\label{52-sm-vs-mssm-coefficients}

{\def\LTcaptype{none} % do not increment counter
\begin{longtable}[]{@{}>{\RaggedRight\arraybackslash}p{0.320\textwidth}>{\RaggedRight\arraybackslash}p{0.320\textwidth}>{\RaggedRight\arraybackslash}p{0.320\textwidth}@{}}
\toprule\noalign{}
Model & \((b_1, b_2, b_3)\) & \(\alpha_s(M_Z)^{\text{pred}}\) \\
\midrule\noalign{}
\endhead
\bottomrule\noalign{}
\endlastfoot
SM & \((41/10, -19/6, -7)\) & \(\approx 0.071\) \\
MSSM & \((33/5, 1, -3)\) & \(\approx 0.116\) \\
Observed &, & \(0.1179 \pm 0.0010\) \\
\end{longtable}
}

The SM prediction is catastrophically wrong; MSSM-like coefficients work.

\subsection{Heat-Kernel Edge Sector Weights}\label{53-heat-kernel-edge-sector-weights}

\textbf{Theorem 5.1:} Under MaxEnt with bi-invariant constraints on a compact Lie group \(G\), the sector probabilities take heat-kernel form: \[p_R(t) \propto d_R \, e^{-t \, C_2(R)}\]

where \(d_R = \dim R\) and \(C_2(R)\) is the quadratic Casimir.

\textbf{Lemma 5.2 (Bi-invariant Operators):} If \(G=\prod_i G_i\) is a compact semisimple Lie group written as a product of compact simple factors, then any bi-invariant, second-order differential operator on \(G\) has the form: \[D = c_0 \mathbf{1} - \sum_i c_i \Delta_{G_i}\]

where \(\Delta_G\) is the Laplace-Beltrami operator.

\emph{Proof:} Bi-invariance \(\Leftrightarrow\) \(D \in Z(U(\mathfrak{g}))\) (center of universal enveloping algebra). For degree \(\leq 2\):

\begin{itemize}
\tightlist
\item
 Linear terms vanish (no invariant vectors in adjoint rep)
\item
 Quadratic terms are proportional to the Killing form on each simple factor, giving one independent coefficient per factor
\end{itemize}

The Casimir element acts as \(-\Delta_G\) in the regular representation. \(\square\)

\subsection{Peter-Weyl and Effective Multiplicity}\label{54-peter-weyl-and-effective-multiplicity}

By Peter-Weyl decomposition, the effective refinement-limit edge representation space is
\[L^2(G) \cong \bigoplus_R V_R \otimes V_R^*\]

\textbf{Key insight:} Entanglement traces over one factor give multiplicity \(d_R\) in \(p_R\), but loops see both factors, giving effective multiplicity: \[N_{\text{eff}}(R) = d_R \cdot p_R\]

\subsection{Beta Function Shifts}\label{55-beta-function-shifts}

The one-loop beta shift from edge sectors: \[\Delta b_a = \sum_R p_R \cdot d_R \cdot T_a(R)\]

where \(T_a(R)\) is the Dynkin index for gauge factor \(a\).

At the unification diffusion parameter \(t_U \approx 1.64\): \[\Delta b \approx (2.49, 4.38, 3.97) \quad \text{vs MSSM} \quad (2.50, 4.17, 4.00)\]

With fermionic grading (restricting to half-integer SU(2) sectors): \[\Delta b \approx (2.50, 4.17, 3.97)\]

matching MSSM within 1\%.

\begin{center}\rule{0.5\linewidth}{0.5pt}\end{center}

\section{The Horizon Problem}\label{6-the-horizon-problem}

\subsection{Anisotropy Bound from Markov Control}\label{61-anisotropy-bound-from-markov-control}

\textbf{Theorem 6.1:} For a tripartition \(A\)-\(B\)-\(D\) with collar width \(\delta\), the Markov error satisfies: \[I(A:D|B) \lesssim c \cdot |\partial C|_{\text{UV}} \cdot e^{-\delta/\xi}\]

where \(\xi\) is the correlation length.

\textbf{Corollary 6.2:} For bounded observable \(O\) with \(\|O\| \sim 1\): \[|\Delta\langle O \rangle| \leq 2\sqrt{\ln 2 \cdot \varepsilon}\]

\subsection{CMB Anisotropy Requirement}\label{62-cmb-anisotropy-requirement}

For \(\delta T/T \lesssim 10^{-5}\): \[2\sqrt{\ln 2 \cdot \varepsilon} \lesssim 10^{-5}\] \[\varepsilon \lesssim 3.61 \times 10^{-11} \text{ bits}\]

Required collar width: \[\frac{\delta}{\xi} \gtrsim \ln\left(\frac{c \cdot |\partial C|_{\text{UV}}}{\varepsilon_{\text{max}}}\right) \approx 24\]

With \(\xi \approx \sqrt{a_{\text{cell}}} \cdot \ell_P \approx 1.28 \ell_P\): \[\delta_{\text{CMB}} \approx 31 \ell_P\]

\subsection{Homogeneity as MaxEnt Default}\label{63-homogeneity-as-maxent-default}

\textbf{Theorem 6.3:} If MaxEnt constraints are uniform across UV cells with no marked points, the MaxEnt state is invariant under triangulation automorphisms. In the continuum limit, this gives SO(3) invariance.

\emph{Proof:} Standard MaxEnt uniqueness + symmetry inheritance (Lemma 4.1). \(\square\)

\begin{center}\rule{0.5\linewidth}{0.5pt}\end{center}

\section{Black Hole Information Paradox}\label{7-black-hole-information-paradox}

\subsection{The Information Trilemma}\label{71-the-information-trilemma}

The AMPS-style paradox assumes:

\begin{enumerate}
\def\labelenumi{\arabic{enumi}.}
\tightlist
\item
 Outgoing mode \(B\) entangled with interior partner \(A\) (smooth horizon)
\item
 Late radiation \(B\) entangled with early radiation \(R\) (unitarity)
\item
 Monogamy: \(B\) cannot be maximally entangled with both
\end{enumerate}

\textbf{The false assumption:} \(\mathcal{H} \stackrel{?}{=} \mathcal{H}_{\text{inside}} \otimes \mathcal{H}_{\text{outside}} \otimes \mathcal{H}_{\text{radiation}}\)

\subsection{Edge-Center Resolution}\label{72-edge-center-resolution}

\textbf{Theorem 7.1:} Under EC, the collar decomposition gives: \[\rho_{A_\delta B_\delta D_\delta} = \bigoplus_\alpha p_\alpha \left(\rho_{A_\delta b_L^\alpha} \otimes \rho_{b_R^\alpha D_\delta}\right)\]

The "glue" between inside and outside is the edge sector label \(\alpha\) in the center. Conditional on \(\alpha\), inside and outside factorize.

\subsection{Recoverability of Interior}\label{73-recoverability-of-interior}

By the recoverability clause of the Recoverable Generalized Entropy axiom (small CMI) + Theorem 3.2: \[\left\|\rho_{ABC} - (\text{id}_A \otimes \mathcal{R})(\rho_{AB})\right\|_1 \leq 2\sqrt{\ln 2 \cdot \varepsilon}\]

\textbf{Interpretation:} The "interior partner" is recoverable from outside + radiation. It\textquotesingle s encoded, not independent.

\subsection{Hawking Temperature from Modular Theory}\label{74-hawking-temperature-from-modular-theory}

\textbf{Theorem 7.2:} The outside algebra is in a KMS state at inverse temperature \(\beta = 2\pi/\kappa\) with respect to the horizon Killing generator.

\emph{Derivation:} Near the horizon: \[ds^2 \approx -\kappa^2 \rho^2 dt^2 + d\rho^2 + r_H^2 d\Omega^2\]

Euclidean regularity at \(\rho = 0\) requires \(\kappa t_E \sim \kappa t_E + 2\pi\), hence: \[T_H = \frac{\kappa}{2\pi} = \frac{\hbar c^3}{8\pi G k_B M}\]

\subsection{Discrete Area Spectrum}\label{75-discrete-area-spectrum}

From the central area operator: \[L_C = \sum_\alpha (\log d_\alpha) P_\alpha\]

Area eigenvalues: \(A_\alpha = 4G \log d_\alpha = 4\ell_P^2 \ln d_\alpha\)

For Schwarzschild, a sector transition \(d \to d'\) gives: \[\Delta E = k_B T_H \ln(d'/d)\]

\begin{center}\rule{0.5\linewidth}{0.5pt}\end{center}

\section{Cosmological Constant and Screen-Capacity Framework}\label{8-the-cosmological-constant-problem}

This section is one local/global package rather than two disconnected remarks. The local null-modular reconstruction fixes the Einstein branch only up to a metric term \(\phi g_{ab}\), so it cannot determine \(\Lambda\) from local overlap data alone. The remaining constant enters only when one supplies the global screen-capacity identification, which then yields the D6 relation \(\Lambda = 3\pi/(G N_{\mathrm{scr}})\).

\subsection{Vacuum Energy Blindness}\label{81-vacuum-energy-blindness}

\textbf{Lemma 8.1:} For any null vector \(k\): \[T^{\text{vac}}_{kk} = T^{\text{vac}}_{ab} k^a k^b = -\rho_{\text{vac}} g_{ab} k^a k^b = -\rho_{\text{vac}}(k \cdot k) = 0\]

Vacuum energy contributions lie in the kernel of the null map \(T_{ab} \mapsto T_{kk}\).

\subsection{Structural Separation}\label{82-structural-separation}

\textbf{Proposition 8.2:} In OPH:

\begin{itemize}
\tightlist
\item
 Local modular/null data fixes \(T_{ab}\) up to \(\phi g_{ab}\)
\item
 \(\Lambda\) is fixed by global screen capacity, not local physics
\end{itemize}

The "huge zero-point energy" simply isn\textquotesingle t the thing determining curvature.

\subsection{\texorpdfstring{\(\Lambda\) from Capacity}{Lambda from Capacity}}\label{83-ux3bb-from-capacity}

\[\Lambda = \frac{3\pi}{G N_{\mathrm{scr}}}\]

Equivalently, using de Sitter entropy: \[S_{dS} = \frac{A_{dS}}{4G} = \frac{3\pi}{G\Lambda} = N_{\mathrm{scr}} = \log(\dim \mathcal{H}_{\text{tot}})\]

For observed \(\Lambda \approx 1.09 \times 10^{-52}\) m\(^{-2}\): \[N_{\mathrm{scr}} \sim 10^{122}\]

Together with Proposition~8.2 and Theorem~8.3, this makes the intended theorematic split explicit: local null data determine the Einstein branch only modulo \(\Lambda g_{ab}\), and the global screen-capacity input fixes the remaining constant.

\subsection{\texorpdfstring{No-Go for Local \(\Lambda\) Prediction}{No-Go for Local Lambda Prediction}}\label{84-no-go-for-local-ux3bb-prediction}

\textbf{Theorem 8.3:} Within OPH\textquotesingle s null-modular reconstruction:

\begin{enumerate}
\def\labelenumi{\arabic{enumi}.}
\tightlist
\item
 All \(T_{kk}\) data unchanged under \(T_{ab} \to T_{ab} + \phi g_{ab}\)
\item
 Therefore \(\Lambda\) cannot be fixed by local overlap consistency
\item
 \(\Lambda\) requires a global input (capacity)
\end{enumerate}

\begin{center}\rule{0.5\linewidth}{0.5pt}\end{center}

\section{Dark-Sector Response from the Modular Anomaly}\label{9-dark-matter-as-modular-anomaly}

Sections 9--13 collect phenomenological branches built on structural inputs proved earlier in the manuscript. They include dark-sector response laws, proton-spin bookkeeping, baryogenesis estimates, and flavor constructions beyond the recovered core. In Section~9, the structural input is the anomalous stress-energy term of Section~9.2 together with the de Sitter scale from D6 once the screen-capacity input is made. Galaxy and cluster phenomenology are not derived in this section.

\subsection{The Modular Additivity Defect}\label{91-the-modular-additivity-defect}

Define the collar modular additivity defect: \[\Delta K_\delta := K_{ABD} - K_{AB} - K_{BD} + K_B\]

\textbf{Theorem 9.1:} \(\langle \Delta K_\delta \rangle_\omega = -I(A:D|B)_\omega\)

This is the conditional mutual information, measuring Markov imperfection.

\subsection{Anomalous Stress-Energy}\label{92-anomalous-stress-energy}

The modified Einstein equation: \[G_{ab} + \Lambda g_{ab} = 8\pi G \left(\langle T_{ab} \rangle + \langle T_{ab}^{\text{anom}} \rangle\right)\]

where in the diamond rest frame: \[\langle T_{00}^{\text{anom}} \rangle := \frac{15}{8\pi^2} \frac{\delta\langle K_C^{(\text{anom})} \rangle}{\ell^4}\]

The coefficient \(15/(8\pi^2)\) comes from inverting the \(d=4\) geometric integral coefficient \(\Omega_2/(4^2-1) = 4\pi/15\).

\subsection{Properties of the Anomaly Sector}\label{93-properties-of-the-anomaly-sector}

\begin{enumerate}
\def\labelenumi{\arabic{enumi}.}
\tightlist
\item
 \textbf{Gravitates:} Appears on RHS of Einstein equation
\item
 \textbf{Dark by construction:} Encoded in central/recoverability data, not SM fields
\item
 \textbf{Covariantly conserved:} \(\nabla^a T_{ab}^{\text{anom}} = 0\) (from Bianchi identity)
\item
 \textbf{Effectively classical:} Central structure implies classical labels
\end{enumerate}

\subsection{Acceleration-Scale Benchmark for the Continuation}\label{94-the-mond-acceleration-scale}

If one tentatively identifies the imported D6 de Sitter scale with the only IR scale available to a modular-anomaly continuation, then a natural acceleration benchmark is \[a_0^{(\text{OPH})} := \frac{15}{8\pi^2} c^2 \sqrt{\frac{\Lambda}{3}} = \frac{15}{8\pi^2} \frac{c^2}{r_{dS}}\]

\textbf{Numerical benchmark:} \[a_0^{(\text{OPH})} \approx 1.03 \times 10^{-10} \text{ m/s}^2\]

This lands near the empirical MOND scale \(a_0^{(\text{obs})} \approx 1.2 \times 10^{-10}\) m/s\(^2\). It does not derive galaxy dynamics, a MOND/RAR law, or an observational dark-matter identification from the recovered core.

\subsection{Illustrative MOND/RAR Ansatz and Missing Inputs}\label{95-polarization-response-mond-like-scaling}

A MOND/RAR-style response such as \[g_{\text{obs}} \approx g_b + \sqrt{a_0 g_b}\]
or the cubic polarization functional \[S_{\text{pol}}[\varphi] = \int d^3x \left[-\frac{1}{12\pi G a_0}|\nabla\varphi|^3 - \rho_b \varphi\right]\]
can be used as phenomenological bookkeeping for the continuation, but neither expression is derived from the OPH axioms or from the D6/D9 stack alone.

What remains missing is a closed dark-sector dynamics package:
\begin{enumerate}
\item a controlled nonrelativistic limit linking \(T_{ab}^{\text{anom}}\) to galaxy observables;
\item a derived source/response law selecting the MOND/RAR functional rather than alternative IR behavior;
\item lensing, cluster, and Bullet-Cluster phenomenology;
\item cosmological abundance and structure-formation checks;
\item environment-dependence and stability bounds.
\end{enumerate}

This modular-anomaly dark-sector branch is a phenomenological continuation rather than a theorem-level closure.

\begin{center}\rule{0.5\linewidth}{0.5pt}\end{center}

\section{Proton Stability and Spin}\label{10-proton-stability-and-spin}

This section separates three distinct claim types:
\begin{enumerate}
\item \textbf{Structural output closed on the realized D8--D9 branch:} the realized low-energy gauge group is a product up to finite quotient, so the adjoint contains no mixed \(X,Y\)-type gauge bosons and therefore no \emph{gauge-mediated} proton-decay channel of simple-GUT type.
\item \textbf{QCD/EFT-dependent hadron claims:} any finite proton-lifetime estimate beyond that gauge-channel exclusion, and any statement about the proton spin decomposition, require explicit operator matching, renormalization conventions, and nonperturbative hadronic matrix elements.
\item \textbf{Completion burden:} the supplement does not contain the nonperturbative light-quark/hadron completion needed to turn those QCD-dependent statements into derived OPH outputs.
\end{enumerate}

\subsection{Sector Factorization and Product Group}\label{101-sector-factorization-and-product-group}

\textbf{Theorem 10.1 (Factorization Equivalence):} \[\text{Factorizing edge weights} \Leftrightarrow \text{Additive boundary Laplacian} \Leftrightarrow \text{Product gauge group}\]

If \(H_\partial = H_\partial^{(1)} + H_\partial^{(2)} + H_\partial^{(3)}\) with \([H_\partial^{(i)}, H_\partial^{(j)}] = 0\): \[p(R_1, R_2, R_3) \propto \prod_{i=1}^3 d_{R_i} e^{-t_i C_2(R_i)}\]

Applied to the refinement-stable edge-sector colimit, Tannaka-Krein reconstruction~\cite{tannaka38,krein49,dr89,dr90} first yields some compact group \(G\). Under the MAR admissibility package, and once the admissible class includes one connected abelian charge factor, the connected Lie realization selected by MAR is \(SU(3) \times SU(2) \times U(1)\) (up to finite quotient).

Under the stated gauge-selection hypotheses, the explicit MAR selector picks this product structure on the admissible class: the minimal coupled carrier \(\mathbb{C}^3 \otimes \mathbb{C}^2\) enforces commuting color and weak actions, which implies the additive boundary Laplacian. See the gauge-branch derivation in the main OPH trunk for the complete proof.

\subsection{No Leptoquarks, No Gauge Proton Decay}\label{102-no-leptoquarks-no-gauge-proton-decay}

\textbf{Corollary 10.2 (structural D8--D9 consequence only):} With product gauge group, the adjoint representation of the full connected gauge group is \[(8,1,0) \oplus (1,3,0) \oplus (1,1,0)\]

equivalently, the adjoint of the derived gauge group is
\[
(8,1,0) \oplus (1,3,0).
\]

There are therefore no gauge generators in mixed representations \((3,2,\pm 5/6)\), i.e.\ no simple-GUT \(X,Y\) bosons on the realized OPH branch. The structural claim is exactly
\[
\tau_p^{(\text{gauge})} = \infty,
\]
meaning that the gauge-boson proton-decay channel is absent.

\textbf{Scope of this result.} The theorem excludes gauge-boson proton decay and does not establish full proton stability or justify the shorthand \(\tau_p=\infty\). Scalar-mediated or higher-dimensional baryon-violating operators, RG transport of those operators from a UV completion, and the hadronic matrix elements entering exclusive proton-decay rates lie outside this structural argument. Those steps are EFT/QCD-dependent. The electroweak baryon-violating counting discussed in Section~\ref{11-baryogenesis} is a separate high-temperature/topological issue, not a zero-temperature proton-stability theorem.

\subsection{Proton Spin Fraction}\label{103-proton-spin-fraction}

\textbf{Structural content fixed.} On the realized D8--D9 branch the color factor is \(\mathrm{SU}(3)\), so the relevant Casimirs are
\[
C_F=\frac43,\qquad C_A=3.
\]
This is the full structural OPH input currently available for proton-spin bookkeeping.

\textbf{QCD-dependent continuation ansatz.} If one adds the phenomenological ansatz that quark/gluon spin sharing equilibrates according to the color Casimirs, then
\[
\Delta\Sigma \approx \frac{C_F}{C_F + C_A}.
\]

For \(\mathrm{SU}(3)\) this gives
\[
\Delta\Sigma = \frac{4/3}{4/3 + 3} = \frac{4}{13} \approx 0.308.
\]

Compared to the representative lattice benchmark \(\Delta\Sigma \approx 0.286\), this lands within about \(8\%\). But the displayed value is only a deferred continuation benchmark, not a theorem. The observable \(\Delta\Sigma\) is factorization-scheme and scale dependent and, in QCD, depends on the polarized axial-current matrix element together with gluon-spin and orbital-angular-momentum contributions.

\textbf{Required inputs.} A first-principles OPH derivation of proton spin would require the nonperturbative light-quark/hadron completion of Phase I plus an explicit map from OPH data to renormalized proton matrix elements. Until that exists, the \(4/13\) estimate is only a QCD-dependent continuation ansatz.

\begin{center}\rule{0.5\linewidth}{0.5pt}\end{center}

\section{Baryogenesis}\label{11-baryogenesis}

This section gives a heuristic bookkeeping branch for baryogenesis. It uses the realized Standard Model quotient and the electroweak topological channel, but it does not derive a dynamical baryogenesis mechanism.

\subsection{Imported structure and missing dynamics}

\begin{itemize}
\tightlist
\item
 \textbf{Imported structural inputs:} the realized Standard Model quotient \((G_{\text{SM}})/\mathbb{Z}_6\), the realized counting chain \(N_g = 3\) then \(N_c = 3\), and the standard electroweak anomalous \(B+L\)-violating channel once the SM matter content is fixed.
\item
 \textbf{Separate upstream point:} no gauge-mediated proton decay is a Phase-I product-group corollary and is \emph{not} the baryon-violating mechanism used in this branch.
\item
 \textbf{Missing dynamical ingredients:} a concrete out-of-equilibrium epoch, a justified defect/sphaleron production dynamics, derived CP-odd source terms, transport and washout control, and a freeze-out computation of the final asymmetry.
\end{itemize}

Until those missing ingredients are derived, any \(\eta_B\) number written below is only a heuristic continuation benchmark.

\subsection{\texorpdfstring{The \(Z_6\) Defect Suppression}{The Z6 Defect Suppression}}\label{111-the-zux2086-defect-suppression}

The SM global gauge structure is \((G_{\text{SM}})/\mathbb{Z}_6\). For this deferred continuation branch we adopt a uniform \(\mathbb{Z}_6\) center-label ensemble as a phenomenological bookkeeping ansatz.
\[
\Delta S = \ln 6.
\]

This ansatz then defines the bookkeeping suppression parameter
\[
\varepsilon = e^{-\Delta S} = \frac{1}{6}.
\]
This is a D12 continuation ansatz, not a D9-derived output, and the OPH package does not prove that a nonequilibrium baryogenesis history weights each relevant event by a fixed power of \(\varepsilon\).

\subsection{The Baryon-Violating Channel}\label{112-the-baryon-violating-channel}

The electroweak topological transition (\(\Delta N_{CS} = \pm 1\)) produces chiral zero modes. By index theorem:

\begin{itemize}
\tightlist
\item
 One zero mode per left-handed SU(2) doublet per unit topological charge
\item
 Per generation: \(N_c + 1 = 4\) doublets
\item
 Total: \(n_{\text{doublets}} = (N_c + 1)N_g = 12\)
\end{itemize}

The \textquotesingle t Hooft vertex is therefore a 12-fermion interaction. This is an imported Standard Model fact once the realized matter package is fixed; it is not itself generated by the \(\mathbb{Z}_6\) bookkeeping parameter.

\subsection{Defect-Order Identification}\label{113-defect-order-identification}

\textbf{Heuristic identification only.} If one assigns one unit of the bookkeeping suppression \(\varepsilon\) per chiral zero mode in the 12-fermion \textquotesingle t Hooft vertex, then the resulting benchmark is
\[
\eta_B^{(\mathrm{heur})} \sim \varepsilon^{12} = 6^{-12} \approx 4.59 \times 10^{-10}.
\]
This is not a theorem. The assignment of one power of \(\varepsilon\) per zero mode is not currently derived, the overall rate prefactor is unknown, and the comparison with observed \(\eta_B \approx 6.1 \times 10^{-10}\) omits transport, washout, and freeze-out. It should therefore be read only as a heuristic continuation benchmark.

\subsection{CP Bias from Overlap Holonomy}\label{114-cp-bias-from-overlap-holonomy}

Under CP: \(U_{ij} \mapsto U_{ij}^*\), \(\Omega_{ijk} \mapsto \Omega_{ijk}^\dagger\)

For central cocycle \(z_{ijk} \in U(1)\): CP sends \(z \to z^{-1}\).

\textbf{Real cocycle condition:} \([z] = [z]^{-1}\) (cohomologous to inverse)

Generic phases violate this \(\Rightarrow\) CP-odd overlap data are generically available in OPH. But this does not derive the CP source terms that would enter a kinetic or Boltzmann treatment. Turning this observation into a baryogenesis theorem requires an effective description that links cocycle phases to defect production, diffusion, washout, and the final baryon asymmetry.

\begin{center}\rule{0.5\linewidth}{0.5pt}\end{center}

\section{Generation Structure and Flavor Continuations}\label{12-generation-structure-and-yukawa-hierarchy}

This section keeps one imported structural corollary and the D12 flavor-status boundary on one surface. The structural corollary is the realized D9 counting chain \(N_g=3\). The D12 advance is narrower: the corpus puts the quark mass-side continuation, the charged same-carrier readback theorem, and the intrinsic neutrino eta-chain on one common excitation dictionary. What remains open are named continuation closures above that dictionary. They are not independent Yukawa parameters, and they do not belong to the recovered-core theorem surface.

\subsection{Imported structural corollary: three generations}\label{121-why-three-generations}

\textbf{Imported D9 corollary (realized counting chain).}

\begin{enumerate}
\def\labelenumi{\arabic{enumi}.}
\tightlist
\item
 \textbf{CP violation requires N\_g \ensuremath{\geq} 3:} Physical CKM phase exists only for \(N_g \geq 3\)
\item
 \textbf{UV stability bounds N\_g \ensuremath{\leq} 5:} SU(2) asymptotic freedom
\item
 \textbf{The realized MAR selector chooses the minimal viable \(N_g\) inside that window:} once the imported D9 branch fixes the viable CP/asymptotic-freedom window, MAR selects \(N_g=3\).
\end{enumerate}

The familiar language that extra generations would introduce unprotected relevant deformations is supporting intuition for why higher-\(N_g\) candidates are disfavored on the realized branch. It is not the corollary statement imported on this surface.

\subsection{Shared excitation dictionary}\label{122-yukawa-as-defect-mediated-overlap}

On every realized same-label arrow \(e\in\{12,23,31\}\), the proof-facing excitation data are the gap and overlap scalars
\[
(g_e,\omega_e),
\qquad
d_e:=1-\omega_e,
\qquad
q_e:=\sqrt{g_e d_e}.
\]
Their centered logarithmic class is
\[
\eta_e:=\log q_e-\frac13\sum_f \log q_f,
\qquad
[\eta_e]\in\mathbb R^3/\langle(1,1,1)\rangle,
\]
equivalently the positive normalized weights
\[
\mu_e:=\frac{e^{\eta_e}}{\frac13\sum_f e^{\eta_f}}.
\]
Let \(s_e\) denote the realized sector-charge step carried by the arrow \(e\).
For any realized lane readout \(\gamma\) on the same-label transport graph, the primitive flavor-side integers are the traversal multiplicities
\[
m_{\gamma,e}\in\mathbb Z_{\ge0}
\]
with which \(\gamma\) traverses the elementary arrows, equivalently the integers satisfying
\[
s_\gamma=\sum_e m_{\gamma,e}s_e .
\]
The associated excitation action is
\[
A_\gamma
:=
-\sum_e m_{\gamma,e}\log q_e .
\]
This common dictionary underlies the quark mass-side continuation, the charged same-carrier readback theorem, and the intrinsic neutrino eta-chain. The older \(6^{-n}\) suppression rule is at most the compressed benchmark \(n_\gamma^{(6)}:=A_\gamma/\log 6\) after an extra uniform-center-label collapse of this richer dictionary and is not a recovered-core theorem.

\subsection{Quark continuation boundary}

On the quark lane, the exact mass-side data are
\[
(\Delta_{ud}^{\mathrm{overlap}},\,\eta_Q^{\mathrm{centered}}),
\]
and on the compact same-family branch they collapse further to the emitted ray
\[
D12_{ud}^{\mathrm{mass}}
\equiv
\mathcal R_{D12}^{ud}.
\]
That ray is the exact same-family mass object emitted by the corpus. What remains open is not another selector search on that surface but the distinguished modulus on the emitted ray. The sole next theorem object is therefore the one-scalar law \ophid{quark_d12_t1_value_law}, which identifies the ray modulus as \(t_1\) on \(D12_{ud}^{\mathrm{mass}}\); until that law is emitted, any numerical scale choice on the ray remains compare-only. The larger object \ophid{intrinsic_scale_law_D12} is retained only as the derived wrapper around that scalar law.

The mixing-side transport
\[
V_{\mathrm{CKM}}^{\mathrm{fwd}}=U_u^\dagger U_d
\]
is also closed on the same selected D12 sheet, but that does not produce a physical branch change. On the emitted same-label left-handed solver surface the selector problem is closed in the negative sense: the local orbit collapses to the singleton \(\sigma_{\mathrm{ref}}\) (\ophid{sigma_ref}), and \(\sigma_{\mathrm{ref}}\) is the wrong D12 CKM sheet. The corpus emits no sector-attached left-handed same-label branch-changing transport from that selected sheet to the physical CKM shell. CKM and CP closure are therefore downstream of the wrong-sheet no-go rather than a second open selector task on the emitted surface. Separately, the builder-facing pure-\(B\) payload pair beneath the reported quark rows remains open.
A target-anchored same-family exact witness reproduces the six quark reference masses exactly on the same ordered three-point family, and within that same-family scope the selected-sheet exact readout closes through \ophid{oph_quark_current_family_selected_sheet_exact_closure}, built on \ophid{oph_quark_current_family_quadratic_readout_theorem}; it does not remove the wrong-branch CKM boundary or emit the one-scalar law.

\subsection{Charged and neutrino continuation boundary}

The charged lane is exact only up to centered same-carrier readback: once a charged source pair exists, the centered charged logs are fixed, but the theorem does not emit absolute masses. The exact sidecar on that surface is the same-family witness \((m_e,m_\mu,m_\tau)=(0.00051099895,0.1056583755,1.7769324651340912)\,\mathrm{GeV}\), and its caveat is exactly that it remains same-family-only. On that same fixed target-anchored witness, the scoped affine coordinate closes through \ophid{oph_lepton_current_family_affine_anchor_theorem}, with \(A_{\mathrm{ch}}^{(\mathrm{sf})}=\tfrac13\log\det(Y_e^{(\mathrm{sf})})\) and \(g_e^{(\mathrm{sf})}=e^{A_{\mathrm{ch}}^{(\mathrm{sf})}}\). The theorem-grade waiting set is promotion of \(\widehat C_e^{\mathrm{cand}}\) and then the post-promotion lift whose descended scalar is \(\mu_{\mathrm{phys}}(Y_e)\); within that lift the exact smaller forcing object is the physical identity-mode equalizer.

The isotropic intrinsic neutrino branch is excluded by the exact atmospheric-splitting cap, so the continuation lane is the weighted-cycle / Majorana-holonomy branch. On that lane the standard-math-fixed balanced transport-load selector is the arithmetic midpoint on the positive segment between \(\chi=1+\varepsilon\) and \(1+\gamma/2\):
\[
D_\nu=\frac{\chi+1+\gamma/2}{2},
\qquad
p_\nu(\gamma,\varepsilon)=1+\gamma+\frac{\varepsilon}{D_\nu}.
\]
This continuation closes the oscillation-shape lane in the physical PMNS/hierarchy regime. It emits
\[
\theta_{12}=34.225904631810025^\circ,\qquad
\theta_{23}=49.72282845058266^\circ,\qquad
\theta_{13}=8.686355527700156^\circ,
\]
\[
\delta_{\mathrm{PMNS}}=305.58061231449796^\circ,\qquad
J=-0.02753115613565372,\qquad
\frac{\Delta m_{21}^2}{\Delta m_{32}^2}=0.030721110097966534.
\]
So no further discrete selector remains on that lane. The normalized same-label overlap-defect weight section closes beneath that shape. The paper-facing amplitude parameterization is
\[
B_\nu:=\lambda_\nu q_{\mathrm{mean}}^{p_\nu}/m_{\star,\mathrm{eV}}
\]
above the closed normalizer, while the best emitted proxy is
\[
P_\nu:=I_\nu^{1/2}\,\widehat{\mathrm{ratio}}^{\,1/2}\,\mathrm{sum\_defect}^{-1}.
\]
Because that proxy is internal to the emitted stack, the smallest exact remaining object is the reduced bridge-correction invariant
\[
C_\nu:=\frac{B_\nu}{P_\nu}.
\]
Equivalently,
\[
B_\nu=P_\nu\,C_\nu,
\qquad
\lambda_\nu=\frac{m_{\star,\mathrm{eV}}}{q_{\mathrm{mean}}^{p_\nu}}\,P_\nu\,C_\nu.
\]
The best closed constructive local object beneath that reduced scalar is the defect-weighted same-label edge family
\[
q_e=\sqrt{g_e d_e},
\qquad
\eta_e=\log q_e-\frac13\sum_f \log q_f,
\qquad
\mu_e=\mu_\nu\,\frac{e^{\eta_e}}{\frac13\sum_f e^{\eta_f}},
\]
but that family does not emit \(C_\nu\) or the induced \(B_\nu\) by itself. The compare-only value is \(C_\nu^\ast=0.9995473560\). A direct correction audit gives the target-containing reduced window \(C_\nu\in[0.9994296,1.0015784]\), which induces the sharpened paper-facing corridor \(B_\nu\in[6.69600,6.71040]\) with midpoint \(6.70320\); the earlier three-route bridge corridor \([6.69315,6.71259]\) remains as a wider comparison envelope, and the shortlist-consensus window \([6.70484,6.71259]\) remains diagnostic-only because it misses the compare-only target. A compare-only two-parameter continuation adapter on the explicit positive selector segment is on disk and hits the representative central splittings exactly, with \((m_1,m_2,m_3)=(0.01745663295,0.01948419960,0.05308139066)\,\mathrm{eV}\), but it remains non-promotable while \(C_\nu\) stays open. So the charged and neutrino lanes remain open for different reasons: charged leptons miss an emitter upgrade and one descended physical affine scalar, while neutrinos miss one reduced positive bridge scalar above the emitted proxy after the shape closure.

\subsection{Compare-only texture benchmarks}\label{123-the-yukawa-spectrum}

The familiar base-6 texture table remains a compare-only continuation benchmark: across the charged-fermion spectrum the quantity
\[
n:=-\frac{\ln y}{\ln 6}
\]
lands near integers, but the integer assignments still depend on extra texture choices, scheme conventions, and running/matching assumptions. That benchmark does not alter the exact open-object ledger above.

\subsection{Optional deterministic toy completion}\label{124-deterministic-completion-rk-hamiltonian}

An RK-type completion can still be written as a toy continuation by assigning a configuration weight
\[
w(\mathcal C)=\exp\!\left(-t\sum_{\text{plaquettes}}C_2(R_p)\right)\times 6^{-n_D(\mathcal C)}
\]
and the corresponding pure state
\[
|\Psi\rangle
=
\frac{1}{\sqrt Z}\sum_{\mathcal C}\sqrt{w(\mathcal C)}\,|\mathcal C\rangle.
\]
This remains an optional continuation ansatz. It is not a proved OPH corollary, and it does not change the sharper quark, charged-lepton, or neutrino frontier stated above.

\begin{center}\rule{0.5\linewidth}{0.5pt}\end{center}
\section{Charged-lepton continuation status}\label{13-the-koide-formula}

This section presents the theorem boundary supported by the corpus. The exact output is the centered same-carrier readback law. The unresolved part is the absolute charged scale, and its remaining obstruction is explicitly two-layered rather than hidden inside a fit summary.

\subsection{Exact centered readback on the ordered charged carrier}\label{131-numerical-verification}

On the ordered charged carrier
\[
(-1,x_2,1),
\qquad
x_2=-0.5175863354681689,
\]
every charged source pair \((\eta,\sigma)\) determines the centered charged logs by
\[
e_{\log,\mathrm{centered}}
=
-\frac{(3+x_2)\sigma-\eta}{6},
\qquad
\mu_{\log,\mathrm{centered}}
=
\frac{x_2\sigma-\eta}{3},
\qquad
\tau_{\log,\mathrm{centered}}
=
\frac{(3-x_2)\sigma+\eta}{6}.
\]
Hence the centered charged hierarchy closes once a charged source pair is emitted. This is the exact theorem-grade charged-lepton output on this surface.

\subsection{The common-shift quotient surface is closed as a no-go}\label{132-zux2083-holonomy-parameterization}

The centered data live on the common-shift quotient: adding \(c\,\mathbf 1\) to an uncentered log vector leaves the centered class unchanged. Therefore no construction that factors only through the centered quotient can emit an absolute charged scale. In particular, the equalizer/common-shift route is closed negatively on the theorem surface: centered data alone do not determine theorem-grade public \(m_e\), \(m_\mu\), or \(m_\tau\).

This is the precise obstruction behind the old absolute-scale gap. The relevant missing object is not another centered invariant and not a Koide-style fitted phase. Any future absolute closure must descend from an uncentered physical object.

\subsection{Remaining theorem-grade frontier}\label{133-why-q--23}

The waiting set is explicitly layered:
\[
\widehat C_e^{\mathrm{cand}}
\longrightarrow
\widehat C_e
\quad\Longrightarrow\quad
\mu_{\mathrm{phys}}(Y_e)
\quad\Longrightarrow\quad
(g_e,\Delta_e^{\mathrm{abs}}).
\]

The first arrow is the package \ophid{oph_generation_bundle_branch_generator_splitting}. Its sharp remaining clause is
\ophid{compression_descendant_commutator_vanishes_or_is_uniformly_quadratic_small_after_central_split}.
The corpus proves neither exact vanishing nor uniform quadratic smallness of that descended commutator after the central split, so \(\widehat C_e^{\mathrm{cand}}\) is not promoted to \(\widehat C_e\). Even a successful closure at this first layer would promote only the centered operator candidate.

The second arrow is not a fresh fitted scalar. The common-shift quotient surface is closed negatively, so centered data alone do not determine an absolute charged scale. The post-promotion slot is the refinement-stable uncentered trace lift \(Y_e\mapsto \widetilde C_e(Y_e)\). Its descended affine scalar is \(\mu_{\mathrm{phys}}(Y_e)\), and the exact smaller forcing object beneath that scalar is \ophid{charged_physical_identity_mode_equalizer}, namely the physical identity-mode equalizer \(\delta(r,r')=0\) on common physical fibers. If theorem-grade \(\widehat C_e\) and refinement-stable physical \(Y_e\) are obtained, then the uncentered lift, determinant-line section, and affine anchor follow canonically:
\[
\widetilde C_e(Y_e)=\widehat C_e(Y_e)+\mu_{\mathrm{phys}}(Y_e)\,\mathbf 1,
\qquad
s_{\det}(Y_e)=3\mu_{\mathrm{phys}}(Y_e),
\qquad
A_{\mathrm{ch}}(Y_e)=\mu_{\mathrm{phys}}(Y_e).
\]
So the coarse affine placeholder \(A_{\mathrm{ch}}\) is not an independent post-promotion datum; it collapses to \(\mu_{\mathrm{phys}}(Y_e)\) on theorem-grade physical \(Y_e\).

\subsection{Same-family exact witness and compare-only continuation boundary}\label{134-the-koide-phase-from-oph}

A target-anchored same-family exact witness reproduces the charged reference triple exactly on the same ordered eigenvalue family, and within that same-family scope the ordered three-point quadratic readout is closed through \ophid{oph_lepton_current_family_quadratic_readout_theorem}. On that same fixed target-anchored witness, the scoped affine coordinate also closes through \ophid{oph_lepton_current_family_affine_anchor_theorem}, with
\[
A_{\mathrm{ch}}^{(\mathrm{sf})}
=
\tfrac13\log\det(Y_e^{(\mathrm{sf})}),
\qquad
g_e^{(\mathrm{sf})}
=
e^{A_{\mathrm{ch}}^{(\mathrm{sf})}}.
\]
This same-family exact witness does not promote either layer of the theorem-grade frontier above.

Balanced Hermitian-circulant, \(Q=2/3\), and \(\delta=2/9\) relations may be used as compare-only continuation surfaces for centered-log audits. They do not promote \(\widehat C_e^{\mathrm{cand}}\), and they do not supply the descended physical affine scalar \(\mu_{\mathrm{phys}}(Y_e)\). The public claim of this section is therefore only the exact centered readback law, the closed common-shift no-go, and the explicit layered remaining frontier.

\section{Cosmological Parameters}\label{14-cosmological-parameters}

\subsection{The de Sitter Scale}\label{141-the-de-sitter-scale}

From \(\Lambda \approx 1.09 \times 10^{-52}\) m\(^{-2}\): \[r_{dS} = \sqrt{\frac{3}{\Lambda}} \approx 1.66 \times 10^{26} \text{ m}\]

The de Sitter time: \[t_\Lambda = \frac{r_{dS}}{c} \approx 5.53 \times 10^{17} \text{ s} \approx 17.5 \text{ Gyr}\]

\subsection{Emergent Age of the Universe}\label{142-emergent-age-of-the-universe}

For flat \ensuremath{\Lambda}CDM: \[t_0 = \frac{2}{3H_0\sqrt{\Omega_\Lambda}} \sinh^{-1}\left(\sqrt{\frac{\Omega_\Lambda}{\Omega_m}}\right)\]

With \(H_0 \approx 67.4\) km/s/Mpc, \(\Omega_\Lambda \approx 0.685\): \[t_0 \approx 13.8 \text{ Gyr}\]

\subsection{Screen Capacity}\label{143-screen-capacity}

\[\dim \mathcal{H}_{\text{tot}} = \exp(S_{dS}) = \exp(N_{\mathrm{scr}}) = \exp\left(\frac{3\pi}{G\Lambda}\right)\]

\[N_{\mathrm{scr}} = \log(\dim \mathcal{H}_{\text{tot}}) \approx 2.85 \times 10^{122}\]

\subsection{Pixel Parameters}\label{144-pixel-parameters}

\begin{itemize}
\tightlist
\item
 Cell area: \(a_{\text{cell}} \approx 1.63 \ell_P^2\)
\item
 UV length: \(\ell_{\text{UV}} = \sqrt{a_{\text{cell}}} \approx 1.28 \ell_P\)
\item
 Entropy per cell: \(\bar{\ell} = a_{\text{cell}}/(4\ell_P^2) \approx 0.408\)
\end{itemize}

\begin{center}\rule{0.5\linewidth}{0.5pt}\end{center}

\section{Appendix A: Summary of Derived Results}\label{appendix-a-summary-of-derived-results}

\subsection{A.1 Theorem-Level Results and Dependency Status}\label{a1-theorem-level-results}

This appendix summarizes the theorem-level results for the main manuscript. The recovered core is D1--D5 together with D7--D9; D6 appears only as the adjacent Phase-II boundary marker. The table keeps imported mathematics, extra premises, and claim tier visible.

\begingroup
\scriptsize
{\def\LTcaptype{none}
\setlength{\tabcolsep}{2pt}
\begin{longtable}[]{@{}>{\RaggedRight\arraybackslash}p{0.07\textwidth}>{\RaggedRight\arraybackslash}p{0.21\textwidth}>{\RaggedRight\arraybackslash}p{0.15\textwidth}>{\RaggedRight\arraybackslash}p{0.22\textwidth}>{\RaggedRight\arraybackslash}p{0.17\textwidth}@{}}
\toprule\noalign{}
Node & Output & Standard mathematics used & Extra premises / inputs beyond the five axioms & Status \\
\midrule\noalign{}
\endhead
\bottomrule\noalign{}
\endlastfoot
\textbf{D1} & schedule-independent normal form for overlap repair & Newman's lemma & finite patch net, termination, and local confluence & Phase I structural theorem \\
\textbf{D2} & EC decomposition, Markov collar, and generalized-entropy split & HJPW; Petz; Fawzi--Renner & explicit recoverability control when exact identities are used; coarse-grained \(A/(4G)\) dictionary & Phase I structural lemma/proposition layer \\
\textbf{D3} & geometric modular flow and Lorentz branch & Bisognano--Wichmann modular geometry; conformal-group classification & T1 together with T3 & Phase I conditional scaling-limit theorem \\
\textbf{D4} & null modular bridge to \(T_{kk}\) & Borchers--Wiesbrock; distributional differentiation on half-lines & the theorem-local inherited-strip decomposition and exact-or-controlled Markov hypotheses of Section~\ref{23-null-modular-bridge-eft-conditions-n1-n3}, plus the downstream density-upgrade and relativistic null-stress-identification premises; the endpoint/propagation controls are internal to the local finite-constraint MaxEnt branch & Phase I conditional bridge theorem \\
\textbf{D5} & Jacobson-type rest-frame relation plus the conditional tensor upgrade & Jacobson small-ball mathematics; null-to-tensor reconstruction modulo a metric term & T4 together with the locally Lorentzian \(d=4\) scaling regime and the all-directions/all-reference-states premise used in the tensor upgrade & Phase I conditional scaling-limit theorem/corollary \\
\textbf{D6} & capacity/\(\Lambda\) corollary and the separate capacity-based neutrino estimate & de Sitter entropy relation and dimensional analysis & external input \(N_{\mathrm{scr}}\) and the screen-capacity identification & Phase II input-dependent corollary / estimate \\
\textbf{D7} & compact gauge reconstruction in the bosonic branch & Doplicher--Roberts / Tannaka reconstruction & T2 and T5--T6 together with T7's directed-colimit/fiber premise for transportable edge sectors; the local-Gibbs/Dobrushin-type collar package supplies only fixed-cutoff support and recoverability control and does not by itself prove existence, persistence, triviality, or nontriviality of the refinement-limit sector colimit & Phase I conditional structural theorem \\
\textbf{D8} & product structure up to finite quotient on the realized branch & compact Lie representation classification; Schur's lemma & MAR, connected admissible class, one connected abelian factor, and vanishing relevant transport obstruction & Phase I realized-branch theorem \\
\textbf{D9} & exact hypercharges, the counting chain \(N_g=3\) then \(N_c=3\), and the \(\mathbb Z_6\) quotient & anomaly algebra; Witten anomaly; CKM CP counting & realized one-generation/one-Higgs package and MAR admissibility premises & Phase I realized-branch theorem/corollary chain \\
\end{longtable}
}
\endgroup

The recovered core in this appendix is D1--D5 together with D7--D9. D6 appears only to keep the Phase-I / Phase-II boundary explicit. D10 and D12 are not theorem-level nodes here.
\subsection{A.2 Secondary Quantitative and Continuation Benchmarks}\label{a2-quantitative-predictions}

This appendix lists the non-core quantitative branches tracked in the manuscript:

\begin{enumerate}
\item \textbf{Phase II input-dependent corollary (D6):} the capacity relation for \(\Lambda\) and the associated order-of-magnitude capacity / neutrino estimate.
\item \textbf{Phase II calibration sector (D10):} gauge-coupling closure of the D10 branch, including the associated beta-shift comparison.
\item \textbf{Phase III deferred continuations (D12 and beyond):} charged-lepton/Koide, texture, neutrino mass/mixing constructions beyond the D6 estimate, dark-sector, baryogenesis, black-hole spectroscopy, string/worldsheet, proton-spin, and any proton-lifetime estimate beyond the structural exclusion of gauge-mediated decay.
\end{enumerate}

These failures localize by tier: item 3 retracts the corresponding continuation only; item 2 challenges D10; item 1 challenges the screen-capacity identification only. None of them erases the recovered core D1--D5 plus D7--D9 by itself.

\subsection{A.3 Open Technical Problems}\label{a3-future-directions}

\begin{enumerate}
\item construct or prove the scaling-limit regime required by the D3--D5 relativity chain, including the null-stress bridge and fixed-cap stationarity premises;
\item extend the compact-gauge reconstruction from the bosonic internal-gauge branch to the fermionic/super-Tannakian chiral branch while preserving the D7--D9 Standard Model structural chain;
\item derive or independently fix the external continuous inputs $P$ and $N_{\mathrm{scr}}$, and complete a fully auditable D10 numerical pipeline with explicit running, threshold, and matching conventions;
\item separate the D6 capacity relation from flavor-style add-ons by deciding exactly which neutrino or cosmological continuations survive without additional ans\"atze;
\item only after items 1--4 are secured, reopen the Phase-III continuations: charged-lepton/Koide structure, texture exponents, dark-sector response laws, baryogenesis, black-hole spectroscopy, string/worldsheet reorganization, proton-spin bookkeeping, proton-lifetime estimates beyond the structural gauge-channel exclusion, and any similar late-stage continuation such as strong-CP proposals.
\end{enumerate}

Interpretive strange-loop / existence narratives lie outside that ledger. This supplement does not define a theorem-usable configuration space for a global emergence map, prove that observer formation and reconstruction close to an internal self-map on such a space, or supply the compactness/contractivity/stability premises needed for an existence or stability theorem.

\begin{center}\rule{0.5\linewidth}{0.5pt}\end{center}

\section{Appendix B: Mixed-Status Numerical Benchmarks}\label{appendix-b-key-numerical-values}

{\def\LTcaptype{none} % do not increment counter
\begin{longtable}[]{@{}>{\RaggedRight\arraybackslash}p{0.240\textwidth}>{\RaggedRight\arraybackslash}p{0.240\textwidth}>{\RaggedRight\arraybackslash}p{0.240\textwidth}>{\RaggedRight\arraybackslash}p{0.240\textwidth}@{}}
\toprule\noalign{}
Quantity / branch & OPH Value & Observed & Status \\
\midrule\noalign{}
\endhead
\bottomrule\noalign{}
\endlastfoot
\(a_0\) (dark-sector continuation) & \(1.03 \times 10^{-10}\) m/s\textsuperscript{2} & \(1.2 \times 10^{-10}\) m/s\textsuperscript{2} & deferred continuation benchmark \\
\(\Delta\Sigma\) (spin continuation; QCD-dependent ansatz) & 0.308 & 0.29 \ensuremath{\pm} 0.03 & deferred continuation benchmark \\
\(\eta_B\) (heuristic baryogenesis benchmark) & \(4.6 \times 10^{-10}\) & \(6.1 \times 10^{-10}\) & deferred continuation benchmark \\
Koide \(Q\) (flavor continuation) & 2/3 & 0.666664 & deferred continuation benchmark \\
Koide \(\delta\) (flavor continuation) & 2/9 & 0.222225 & deferred continuation benchmark \\
\(\Delta b_3/\Delta b_2\) (D10 calibration) & 0.91 & 0.96 (MSSM) & supplement-backed calibration \\
\end{longtable}
}

\begin{center}\rule{0.5\linewidth}{0.5pt}\end{center}

\section{References}\label{references-technical-supplement}

The derivations in this supplement are based on the OPH framework as specified in the main OPH theorem ledger. Standard results from:

\begin{itemize}
\tightlist
\item
 Tomita-Takesaki modular theory
\item
 Quantum Markov chain structure theorems (Hayden-Jozsa-Petz-Winter)
\item
 Fawzi-Renner recovery bounds
\item
 Gleason\textquotesingle s theorem
\item
 Jacobson\textquotesingle s entanglement equilibrium
\item
 Peter-Weyl decomposition
\item
 Heat kernel asymptotics on compact Lie groups
\item
 Tannaka-Krein reconstruction
\item
 \textquotesingle t Hooft anomaly matching and instanton calculus
\end{itemize}
